\documentclass[11pt]{article}
\usepackage{../../../tex/lecturestyle}

\title{DP Acceleration via Eigenfunctions}
\author{}
\date{}

\begin{document}
\thispagestyle{fancy}

\lecturetitleband{Polynomials \& Formal Power Series}{DP Acceleration via Eigenfunctions}{9}

\section{Overview}
Matrix exponentiation is a familiar trick: to apply a linear map $M$
times, diagonalize it (find its eigenvalues and eigenvectors), and
raising a diagonal matrix to the $M$-th power is just raising each
eigenvalue to the $M$-th power. This lecture applies the same idea to a
different setting: DP recurrences $A_{i+1} = f(A_i)$ where each $A_i$ is
a \emph{polynomial} (or formal power series) and $f$ is a linear map
involving differentiation. When $f$'s eigenfunctions are known in closed
form, $M$ iterations of the DP again collapse to one change of basis, one
exponentiation of each eigenvalue, and one change back. Assumes linear
algebra (eigenvalues/eigenvectors), formal power series, and (for two of
the three worked problems) power series composition and separable
differential equations.

\section{Eigenvalues and Eigenfunctions}
\begin{definitionbox}
For a vector space $V$ and linear map $f: V \to V$, a nonzero
$v \in V$ and scalar $\lambda$ with $f(v) = \lambda v$ make $\lambda$ an
\emph{eigenvalue} and $v$ an \emph{eigenvector} (or \emph{eigenfunction},
when $V$'s elements are themselves functions).
\end{definitionbox}

\begin{examplebox}
For $A = \begin{pmatrix}2&1\\1&2\end{pmatrix}$ and $f(x)=Ax$: the vector
$(1,1)$ is an eigenvector with eigenvalue $3$, and $(1,-1)$ is an
eigenvector with eigenvalue $1$. Separately, if $V$ is the space of all
formal power series and $f$ is differentiation, then for any constant
$\lambda$, $e^{\lambda x}$ is an eigenfunction with eigenvalue $\lambda$
(since $\frac{d}{dx}e^{\lambda x} = \lambda e^{\lambda x}$).
\end{examplebox}

\begin{theorembox}
If $v_0,\ldots,v_n$ are eigenvectors of $f$ with eigenvalues
$\lambda_0,\ldots,\lambda_n$, and $v = \sum_k c_k v_k$, then
\[
  f^M(v) \;=\; \sum_{k=0}^n c_k\, \lambda_k^M\, v_k.
\]
\end{theorembox}
This is the whole mechanism: express the DP's initial state in the
eigenbasis, raise each eigenvalue to the $M$-th power (fast, via binary
exponentiation or a closed form), and convert back.

% FIGURE: eig_square
\begin{figure}[H]
\centering
\begin{tikzpicture}[node distance=2.6cm and 5.2cm]
\node[lnode,minimum width=3.4cm] (A0) {$A_0$ (initial polynomial)};
\node[lnode,minimum width=3.4cm,right=of A0] (AM) {$A_M=f^{M}(A_0)$};
\node[lnode,minimum width=3.4cm,lhi,below=of A0] (C) {coefficients $c_0,\dots,c_N$\\ in the eigenbasis};
\node[lnode,minimum width=3.4cm,lhi,below=of AM] (CM) {$c_0\lambda_0^M,\dots,c_N\lambda_N^M$};
\draw[larrow,lectureink!45,dashed] (A0) -- (AM) node[midway,above,llabel] {$M$ DP steps ($O(M)$ iterations)};
\draw[larrow] (A0) -- (C) node[midway,left,llabel,align=right] {expand in the\\eigenbasis};
\draw[larrow] (C) -- (CM) node[midway,above,llabel] {multiply $c_k$ by $\lambda_k^M$};
\draw[larrow] (CM) -- (AM) node[midway,right,llabel,align=left] {recombine};
\end{tikzpicture}
\caption{The eigenfunction method as a commuting square. Instead of applying the linear map $f$ a total of $M$ times (dashed), change to a basis in which $f$ is diagonal, raise each eigenvalue to the $M$-th power, and change back. The two basis changes cost the same as one Taylor shift (or one composition) each.}\label{fig:eig_square}
\end{figure}

\section{Worked Problems}
The technique is best understood through examples of increasing
difficulty.

\subsection{A tractable case: \texorpdfstring{$A_{i+1}=(1+x)A_i'$}{A\_{i+1}=(1+x)A\_i'}}
\begin{definitionbox}
Given a polynomial $A_0(x)$ of degree $\le N$, define
$A_{i+1}(x) = (1+x)A_i'(x)$. Compute $A_M(x) \bmod 998244353$.
\end{definitionbox}

Here $f(A) = (1+x)A'$ is linear on the $(N{+}1)$-dimensional space of
degree-$\le N$ polynomials. Finding its eigenfunctions means solving
$(1+x)A'(x) = \lambda A(x)$ --- a separable ODE. Writing $y = A(x)$ and
formally separating variables,
\[
  \frac{y'}{y} = \frac{\lambda}{1+x}
  \;\Longrightarrow\;
  (\log y)' = \bigl(\lambda \log(1+x)\bigr)'
  \;\Longrightarrow\;
  y = (1+x)^\lambda .
\]
(The manipulation is purely heuristic --- log of a formal power series,
division by $y$ --- but the conclusion is easy to check directly:
$(1+x)\cdot\lambda(1+x)^{\lambda-1} = \lambda(1+x)^\lambda$, confirming
$(1+x)^\lambda$ really is an eigenfunction with eigenvalue $\lambda$,
regardless of how it was found.) Restricting to degree $\le N$
polynomials, the eigenfunctions are $(1+x)^k$ for $k=0,\ldots,N$, with
eigenvalue $k$.

\begin{theorembox}
Problem 2 above is solved by: (1) find $c_0,\ldots,c_N$ with
$A_0(x) = \sum_k c_k(1+x)^k$; (2) replace each $c_k$ with $c_k \cdot
k^M$; (3) output $\sum_k c_k (1+x)^k$. Both basis-change steps are a
\emph{Taylor shift} (substituting $x \to x - 1$ and back), computable in
$O(N \log N)$ with a fast implementation (the reference code below uses
a simpler $O(N^2)$ version). Total: $O(N(\log N + \log M))$.
\end{theorembox}

\begin{notebox}
Why Taylor shift finds the $c_k$: setting $y = 1+x$, the expansion
$A_0(x) = \sum_k c_k(1+x)^k$ is exactly $A_0(y-1)$ written as a
polynomial in $y$ --- i.e.\ $c$ is $A_0$'s coefficient list after
substituting $x \to x-1$. Converting back at the end is the same
substitution in reverse, $x \to x+1$.
\end{notebox}

\subsection{A subtler case: \texorpdfstring{$A_{i+1}=A_i+A_i'$}{A\_{i+1}=A\_i+A\_i'}}
\begin{definitionbox}
Given $A_0(x)$ of degree $\le N$, define $A_{i+1}(x) = A_i(x) + A_i'(x)$.
Compute $A_M(x) \bmod 998244353$.
\end{definitionbox}

The eigenfunction here is easy to guess directly: $\frac{d}{dx}e^{kx} =
k e^{kx}$, so $e^{kx}$ satisfies $A + A' = (k+1)A$, eigenvalue $k+1$. The
complication is truncation: $e^{kx} \bmod x^{N+1}$ is \emph{not} exactly
an eigenfunction of the corresponding truncated map, because computing
its derivative from the truncated series silently drops information that
the true (infinite) series would have supplied at the top degree ---
verified concretely in the reference code below, where the identity
$A+A' = (k+1)A$ holds at every coefficient except deliberately the
highest one once the working precision is cut too short.

The fix is to work at higher precision than needed and truncate only at
the very end: expand $A_0(x) \equiv \sum_{k=0}^{N+M} c_k e^{kx}
\pmod{x^{N+M+1}}$ (enough headroom that the top-degree error never
reaches the part of the answer that matters), multiply each $c_k$ by
$(k+1)^M$, recombine, and \emph{then} truncate to degree $N$. Finding the
$c_k$ (expanding in the $e^{kx}$ basis) is itself a formal power series
composition problem --- writing $A_0$ as a polynomial evaluated at
$e^x$ --- solvable with a general composition algorithm, giving
$O((N+M)\log^2(N+M))$ overall.

\subsection{A generalized case: \texorpdfstring{$A_{i+1}=x(A_i+A_i')+b_iA_i$}{A\_{i+1}=x(A\_i+A\_i')+b\_i A\_i}}
\begin{definitionbox}
Given $A_0(x)$ of degree $\le N$ and $b_0,\ldots,b_{M-1}$, define
$A_{i+1}(x) = x(A_i(x)+A_i'(x)) + b_i A_i(x)$. Compute $A_M(x) \bmod
998244353$. (This recurrence appears in one solution to AtCoder ABC272
Ex.)
\end{definitionbox}

Direct substitution confirms $x^k e^{-x}$ is an eigenfunction of
$A \mapsto x(A+A')$ with eigenvalue $k$: writing $A = x^k e^{-x}$,
$A' = kx^{k-1}e^{-x} - x^k e^{-x}$, so
$A + A' = kx^{k-1}e^{-x}$, and $x(A+A') = k\,x^k e^{-x} = kA$.
Since the $b_i$ term only rescales by a constant at each step (not mixing
different eigenfunctions), the effective eigenvalue across all $M$ steps
becomes the \emph{product} $(k+b_0)(k+b_1)\cdots(k+b_{M-1})$ rather than a
single power $k^M$.

\begin{theorembox}
The strategy: (1) expand $A_0(x) \equiv \sum_{k=0}^N c_k\, x^k e^{-x}
\pmod{x^{N+1}}$ (multiplying by $e^x$ turns this into an ordinary
polynomial-coefficient read-off, no composition needed here since the
basis functions are already indexed by monomials $x^k$ times a fixed
$e^{-x}$); (2) replace each $c_k$ with $c_k \cdot (k+b_0)\cdots
(k+b_{M-1})$, computable for all $k=0,\ldots,N$ at once via multipoint
evaluation of $(x+b_0)\cdots(x+b_{M-1})$; (3) recombine and multiply back
by $e^{-x}$. Total: $O((N+M)\log^2(N+M))$.
\end{theorembox}

\section{Complexity Summary}
\begin{center}
\begin{tabular}{ll}
\toprule
\textbf{Problem} & \textbf{Complexity} \\
\midrule
$A_{i+1}=(1+x)A_i'$ (eigenbasis $(1+x)^k$) & $O(N(\log N + \log M))$ \\
$A_{i+1}=A_i+A_i'$ (eigenbasis $e^{kx}$) & $O((N+M)\log^2(N+M))$ \\
$A_{i+1}=x(A_i+A_i')+b_iA_i$ (eigenbasis $x^ke^{-x}$) & $O((N+M)\log^2(N+M))$ \\
\bottomrule
\end{tabular}
\end{center}

\section{Reference Implementations}
Both examples below fully implement and test Problem 2 (via Taylor
shift, cross-checked against a direct $O(NM)$ simulation of the DP
table). For Problems 3 and 4, which need a full power-series composition
or multipoint-evaluation pipeline to reach the stated complexity, the
examples instead numerically verify the underlying eigenfunction
identities themselves --- via a from-scratch power series $\exp$ ---
including a concrete demonstration of the truncation caveat described
above.

\codenote
\subsection{Rust}
\corelisting{Rust}{example.rs}

\subsection{Python}
\corelisting{Python}{example.py}

\section{Worked Example}
\begin{examplebox}
For Problem 2, the eigenfunction method (Taylor shift by $-1$, scale by
$k^M$, Taylor shift by $+1$) matches a direct $O(NM)$ simulation of the
stated DP exactly across 200 random $(A_0, M)$ pairs in both reference
implementations. Separately, computing $e^{kx} \bmod x^{10}$ from scratch
and checking $A+A'=(k+1)A$ confirms the identity holds exactly with a few
extra digits of precision kept, and \emph{specifically fails at the
highest computed coefficient} --- and nowhere else --- when that
headroom is removed, exactly matching the lecture's own description of
why the $N{+}M$ padding trick is necessary.
\end{examplebox}

\section{Further Reading}
\begin{itemize}[leftmargin=1.5em]
  \item AtCoder Algorithm Lectures index (source of this curriculum):
        \url{https://info.atcoder.jp/entry/algorithm_lectures/index}
  \item AtCoder ABC272 Ex, ``Flipping Coins 2'' (an application of the
        third worked problem):
        \url{https://atcoder.jp/contests/abc272/tasks/abc272_h}
\end{itemize}

\end{document}
